\talkslide{advprompter}{AdvPrompter: predicting adversarial prompts}{Timing: provisional}{Use an LLM to predict a jailbreak from the initial prompt, amortizing the search effort across prompts.}{
\advprompterflow{0}
\bottompaperrefs{advprompter,amortization}
}
\speak{Previously, jailbreaking even a single prompt was considered extremely difficult, and substantial search computation went into solving each instance. Our insight was to learn across those instances: an LLM can predict a jailbreak from the initial prompt. More precisely, AdvPrompter predicts an adversarial suffix to append to the prompt. This is amortized optimization: the cost of solving attacks during training is spread across many prompts, so generating an attack for a new prompt can reuse what the model has learned. My tutorial gives the general perspective; I would mention it briefly here rather than introduce its notation. AdvPrompter learns an input-conditioned attack, rather than one universal suffix. It trains against a fixed target, alternating search for stronger suffixes with supervised updates to the attacker. It was an early learned attack generator, not the first learned red-teaming method; AmpleGCG appeared contemporaneously in April 2024. Sources: AdvPrompter, Introduction and Section 3.1, Remark on relation to amortized optimization; Brandon Amos, Tutorial on Amortized Optimization, Foundations and Trends in Machine Learning 2023, Chapter 2. The preprint of AdvPrompter appeared in 2024 and the publication is ICML 2025.}
